\runninghead{Shew Us Thy Mercy}
\properlane{mercy-shown-salvation-given}{Shew Us Thy Mercy: Christ Given as Man, to Be Given as God}

Twice the Mass sings from Psalm 84 (85). Before the Gospel it asks,
\latin{Osténde nobis, Dómine, misericórdiam tuam: et salutáre tuum da nobis},
``Shew us, O Lord, thy mercy; and grant us thy salvation'' (v.~8). At
Communion it promises, \latin{Dóminus dabit benignitátem: et terra nostra
dabit fructum suum}, ``The Lord will give goodness: and our earth shall yield
her fruit'' (v.~13). The first verse raises a question for Christians who sing
it. If the salvation God was asked to give is Christ, and Christ has been
given, why does the Church still ask for him? The second raises a question of
its own: what is the earth, and what is its fruit? This interpretation follows
St Augustine's sermon on the psalm and St Robert Bellarmine's explanation of
it.

\subsection{``This is what we sang''}

Augustine preached on Psalm 84 to a congregation that had sung it. When
he reaches v.~8 he says so: \latin{Hoc est quod cantavimus}, ``This is what we
sang'' (\work{Enarrationes in Psalmos} 84.9). He reads the salvation as a
person: \latin{salutare tuum, Christum tuum}, ``thy salvation, thy Christ.''
And he says why Christians still sing it. \latin{Jam quidem
dedit nobis Christum suum: adhuc illi tamen dicamus, Da nobis Christum tuum},
``He has already given us his Christ; yet let us still say to him, Give us thy
Christ,'' as we say to him, \latin{Panem nostrum quotidianum da nobis hodie},
``Give us today our daily bread''; and the bread is the one who said, ``I am
the living bread which came down from heaven'' (Jn 6:51).

Then he explains why the prayer is still needed. \latin{Dedit enim nobis
Christum, sed hominem: quem nobis dedit hominem, eum ipsum nobis daturus est
Deum}, ``For he gave us Christ, but as man: the same whom he gave us as man he
will give us as God.'' He gave him as man because men could receive him only
so; no human being could grasp Christ as God. The second gift is the sight of
God, and Augustine explains it by the adoption of sons: \latin{Jam quidem
sumus, sed per fidem: sumus quidem in spe, nondum sumus in re}, ``We are so
already, but by faith; we are so in hope, not yet in reality.'' He quotes St
John: ``Dearly beloved, we are now the sons of God: and it hath not yet
appeared what we shall be. We know that when he shall appear we shall be like
to him: because we shall see him as he is'' (1~Jn 3:2).

To make his hearers feel what that promise means, Augustine tells a story. A
poor man wandering abroad, who does not know his family and lives in want and
labour, is told one day: ``You are a senator's son; your father enjoys a great
estate; I am calling you back to him.'' How he would rejoice, if the one who
said it were not a deceiver. St John, an apostle of Christ and no deceiver,
says to those who despair, ``You have a father, you have a homeland, you have an
inheritance.'' Augustine ends the section with the prayer the verse makes:
``Give us thy Christ, let us know thy Christ, let us see thy Christ; not as the
Jews saw him and crucified him, but as the angels see him and rejoice''
(84.9).

The mercy shown in the verse's first half has its own sense for Augustine.
God shows his mercy to a man by persuading him ``that whatever good he has, he
has only from him who is all our good.'' Seeing that, a man is not proud, and
Augustine draws out the chain: not being proud he is not lifted up, not lifted
up he does not fall, not falling he stands, standing he clings, clinging he
remains, remaining he enjoys the Lord his God (84.9).

\subsection{Bellarmine: the prayer for the Saviour's coming}

Bellarmine reads the whole psalm as a prophecy of Christ; in John
O'Sullivan's English it is headed ``The Coming of Christ to Bring Peace and
Salvation to Man.'' In the first verses
God is thanked for a restoration decreed; then the prophet prays for its
execution. On v.~8 he writes: ``Having asked that the divine wrath may be
mitigated; and having asked for that reconciliation and regeneration that
always accompanies remission of sin, he now asks for the coming of the
Savior.'' The words ``grant us thy salvation'' mean ``Send us your Son for a
Savior, for then you will clearly show unto all the extent of your mercy''
(\work{Explanatio in Psalmos}, Ps 84, on v.~8).

Augustine and Bellarmine agree that the salvation of v.~8 is Christ given in
the flesh. Augustine says God \latin{dedit \dots\ hominem}, ``gave him as
man''; Bellarmine's prophet asks, ``Send us your Son.'' Augustine looks
further, to the same Christ to be given as God, and he explains that second
gift by the sight promised in St John's letter.

St Thomas Aquinas names the Epistle's salvation in the same way when he
explains how the Church hears it in Advent. Read in this season, he says,
Paul's ``now our salvation is nearer than when we believed'' refers to the
salvation Christ brought at his first coming, and Paul speaks for all the
faithful since the beginning of the world. The Church, in
Advent, represents the time when the Incarnation was drawing near and the
prophecies were growing more frequent, and in that time one could say,
\latin{nunc propior est nostra salus, scilicet Christus}, ``now our
salvation, that is, Christ, is nearer'' (\work{Super Epistolam ad Romanos}
13, lect.~3). If in Advent the Church takes the place, as Aquinas says, of the
time when the Incarnation was drawing near, she stands where Bellarmine puts
the psalmist, among those who ask for the Saviour before he comes, and the
Alleluia is her prayer: ``Send us your Son.''

\subsection{Two gifts and the Mass's two horizons}

The Postcommunion asks to be prepared for \latin{reparatiónis nostræ ventúra
solémnia}, the coming feast of our restoration. The Gospel describes the Son
of man ``coming in a cloud, with great power and majesty.'' Augustine's two
gifts fall on those two horizons of the Mass: Christ given as man, whose birth
the Church prepares to keep, and Christ to be given as God, when ``we shall see
him as he is.'' The two comings of the Son of God that Augustine names, on
Psalm 24, are a different pair: ``the one in mercy, the other in judgment''
(\work{En.\ in Ps.} 24.10).

St Cyril of Jerusalem teaches the two comings in plain terms: ``We preach not
one advent only of Christ, but a second also, far more glorious than the
former.'' The first was ``the unobserved, like rain on a fleece''; in it ``He
was wrapped in swaddling clothes in the manger; in His second, He covereth
Himself with light as with a garment'' (\work{Catechetical Lectures} 15.1).
St Gregory the Great draws the same contrast from the Gospel's words ``with
great power and majesty'': those who would not hear Christ in his humility
will see him in power (\work{Homiliae in Evangelia} I.1.2).

\subsection{The difficulty: what the earth yields}

The commentators disagree about what the Communion's earth and fruit are.

Rupert of Deutz, expounding this chant in his commentary on the Mass, reads
it of the Incarnation: \latin{terra nostra, id est, uterus Virginis
concipiens, dedit fructum suum, Deum et hominem proferens}, ``our earth, that
is, the womb of the Virgin conceiving, gave its fruit, bringing forth God and
man'' (\work{De divinis officiis} III.2). Sicard of Cremona repeats the
reading (\work{Mitrale} V). Bellarmine gives it its fullest
form. The psalm,
he says, ``still treats of the mystery of the Incarnation, showing that truth
could spring out of the earth; not in the manner of the seed that we sow and
cultivate, but in the manner of the natural flowers that grow spontaneously,
with no other culture than the beams of the sun, and the rains of heaven.''
The Lord's goodness is the Holy Spirit, ``who will overshadow a virgin, and
thus our land, which was never ploughed nor sown, and was altogether an
untouched virgin, will yield her fruit'' (\work{Explanatio}, Ps 84, on v.~13).

Augustine reads the same verse of grace in the believer. His psalter had
\latin{Dominus dabit suavitatem}, ``the Lord will give sweetness.'' He takes
the verse before it, ``Truth is sprung out of the earth,'' of the confession
of sins, and ``justice hath looked down from heaven'' of the justification
God gives to the one who confesses. The sweetness God gives is
\latin{suavitatem justitiae faciendae}, ``the sweetness of doing justice,'' so
that a man who delighted in drunkenness delights in sobriety, a man who stole
looks for someone to give to, and \latin{currere ad ecclesiam, qui primo
currebas ad theatrum}, ``you run to the church, who used to run to the
theatre.'' Then he turns to his hearers as a field. He has scattered the seed
of the word, finding their hearts furrowed by ``the plough of confession'';
let them break the clods so that the birds do not take the seed. \latin{Et
nisi Deus pluerit, quid prodest quod seminatur?} ``And unless God sends rain,
what use is what is sown?'' \latin{Veniat imber Dei}, ``Let God's rain come''
(84.15). Sicard gives a second reading, close to Augustine's, for the
churches that read this Mass's Gospel, \latin{Erunt signa}: there the
goodness is a good affection, the earth is our flesh, and its fruit is good
works. Durandus copies him (\work{Mitrale} V; \work{Rationale} VI).

For Rupert and Bellarmine the Communion sings of Christ born of the Virgin,
the fruit the earth gave once. For Augustine it sings of the good works God's
sweetness makes grow in those who confess, a fruit the earth yields again
wherever his rain falls. Both readings agree that the earth bears nothing of
itself: Bellarmine's land was ``never ploughed nor sown,'' and Augustine's
field needs God's rain. The first sense suits a Mass that prepares for
Christmas; the second suits the moment of Communion, when the faithful
receive what they have asked for.

\subsection{Christ asked for, waited for and received}

The Alleluia asks the Father to give his Christ; the Collect, addressed to the
Son, asks him directly, \latin{Excita, quǽsumus, Dómine, poténtiam tuam, et
veni}, ``Bestir, O Lord, Thy might, we pray Thee and come.'' Both prayers ask
for one who has already been given, and Augustine's reason for still singing
the first, that Christ given as man is still to be given as God, serves for
the second.

The Introit, the Gradual and the Offertory sing that none who wait for God are
confounded, and the Offertory sings it as the gifts are brought to the altar.
Bellarmine explains such waiting as expectation in the certain hope of help.
Since the Alleluia and the Collect ask for Christ, the one waited for in this
Mass may be understood as Christ. The Secret then asks that the holy gifts
bring us purer \latin{ad suum princípium}, to their source, and Schuster reads the
request as a prayer that ``we may worthily receive him who is its Author''
(p.~322). The Author received would be the Christ whom the Alleluia asked
God to give.

The noun \latin{misericórdia} appears twice in the Mass. In the Alleluia the
Church asks God to show it: \latin{Osténde nobis \dots\ misericórdiam tuam}.
After Communion she asks to receive it: \latin{Suscipiámus, Dómine,
misericórdiam tuam in médio templi tui}. The Postcommunion borrows the words
of Psalm 47, \latin{Suscepimus}, ``We have received thy mercy,'' and turns
the psalm's thanks into a request. On v.~8 Augustine took the mercy God shows
to be the persuasion that whatever good a man has is from God. With the
same word in both prayers, that is also what the communicant would be asking
to receive, and the prayer asks it so that we may go before the coming feast
of our restoration with fitting honours.

\begin{fourSenses}
\item[Literal.] The psalmist, who has thanked God for turning away his anger,
asks him to show his mercy and give his salvation, and hears the promise that
the Lord will give goodness and the land yield its harvest (Ps 84:8, 13).
\item[Allegorical.] The salvation is Christ (Augustine; Bellarmine), given to
us as man and to be given as God (Augustine). The earth that yields its fruit
is the Virgin, unploughed and unsown, who bears Christ by the overshadowing
of the Spirit (Bellarmine; Rupert).
\item[Moral.] God shows his mercy by teaching a man that all his good is from
God, and so keeps him from pride (Augustine). The earth that confesses its
sins receives the sweetness of doing justice and yields good works, if God's
rain falls on the seed of the word (Augustine).
\item[Anagogical.] We are sons of God now by faith and in hope; when Christ
appears we shall be like him, because we shall see him as he is, as the angels
see him and rejoice (Augustine, with 1~Jn 3:2).
\end{fourSenses}
